% ============================================================
%  CAPÍTULO 7: DISCUSIÓN
% ============================================================
\chapter{Discusión}
\label{cap:discusion}

% ------------------------------------------------------------
\section{Introducción}
% ------------------------------------------------------------

Los resultados presentados en el Capítulo~\ref{cap:resultados} son estadísticamente
robustos y teóricamente consistentes. Sin embargo, su interpretación requiere
contextualización en el debate académico sobre la naturaleza económica del seguro
sanitario privado y las implicaciones de política sanitaria derivadas. Esta sección
articula los hallazgos en torno a tres ejes: (1) la clasificación del seguro como
bien de lujo y sus implicaciones de equidad, (2) el papel de las listas de espera
como mecanismo de segmentación del sistema sanitario, y (3) las implicaciones para
la política sanitaria española.

% ------------------------------------------------------------
\section{El seguro privado como bien de lujo: implicaciones de equidad}
% ------------------------------------------------------------

La elasticidad renta estimada ---1{,}02 según el modelo de efectos fijos
preferido por el test de Hausman ($\chi^2(4) = 11{,}44$, $p = 0{,}022$)--- sitúa el
seguro sanitario privado en España en la frontera entre \textbf{bien normal y
bien de lujo}, coherente con los resultados de estudios previos para otros países
industrializados. Bajo efectos aleatorios la elasticidad asciende a 1{,}54, lo
que reforzaría la clasificación como bien de lujo.

Esta evidencia tiene consecuencias fundamentales para la equidad del sistema
sanitario español. La hipótesis de que la sanidad es un bien de mérito ---cuyo
consumo debe ser independiente de la renta--- es cuestionada por la evidencia: los
hogares de mayor renta concentran el acceso al seguro privado, que proporciona
tiempos de espera menores, mayor acceso a especialistas y mayor control sobre el
momento de la intervención quirúrgica.

La elasticidad superior a la unidad implica que, ante un crecimiento económico,
la proporción del gasto sanitario privado en la renta aumenta, lo que puede
traducirse en una \textbf{segmentación creciente} del sistema sanitario. Este
patrón es compatible con lo que \textcite{Hirschman1970} denominó el efecto
\textit{exit}: los grupos de mayor renta abandonan progresivamente el sistema
público, reduciendo su base imponible política y debilitando su dotación de
recursos.

% ------------------------------------------------------------
\section{La espera como precio sombra: mecanismo de segmentación}
% ------------------------------------------------------------

El análisis de los retardos de las listas de espera (M2-FE) muestra que los dos
retardos no son conjuntamente significativos ($\chi^2(2) = 0{,}39$, $p = 0{,}82$).
Esto sugiere que el efecto lineal de los días de espera no opera de forma
acumulativa a través del tiempo.

Sin embargo, el \textbf{efecto umbral resulta significativo}: cuando la espera
supera los 100 días, el gasto en seguros privados aumenta un \textbf{9{,}3\%}
(M1-U: $\hat{\delta} = 0{,}093$, $p = 0{,}017$). Este resultado es central para
la investigación: el efecto de las listas de espera \textbf{no es lineal}, sino
que opera como un mecanismo de ``disparo'' cuando se supera un umbral crítico.
Esto es coherente con el efecto \textit{exit} de Hirschman: los hogares no ajustan
gradualmente su comportamiento ante pequeños incrementos de espera, sino que
``saltan'' al sector privado cuando perciben que el sistema público ha cruzado
un límite de tolerancia.

Además, el modelo M3-FE revela una \textbf{persistencia del gasto} al borde de
la significatividad: el logaritmo del gasto del año anterior predice el gasto
actual con una elasticidad de $\hat{\rho} = 0{,}210$ ($p = 0{,}075$), lo que
refleja la inercia en las decisiones de contratación de seguros ---probablemente
por renovación automática de pólizas y formación de hábitos de consumo
sanitario privado.

% ------------------------------------------------------------
\section{Implicaciones para la política sanitaria española}
% ------------------------------------------------------------

\subsection{La función de la renta en el diseño de la política sanitaria}

La elevada elasticidad renta implica que las condiciones macroeconómicas tienen
un impacto directo sobre la composición del sistema sanitario. En períodos de
crecimiento económico, el sector privado crece más que proporcionalmente, lo que
puede aliviar temporalmente la presión sobre el sistema público pero refuerza
la estratificación del acceso.

La política sanitaria debe tener en cuenta que, en ausencia de intervención,
el crecimiento económico tiende de forma natural a ampliar las
\textbf{brechas de acceso} entre ricos y pobres dentro del sistema sanitario.
Esto no ocurre porque el sistema público empeore en términos absolutos, sino
porque el privado mejora más rápidamente en calidad percibida por los hogares de
alta renta.

\subsection{La gestión de las listas de espera como política de cohesión}

El resultado de que el efecto umbral es significativo (H4 confirmada) implica que
las listas de espera tienen un \textbf{doble coste social}: el directo (morbilidad
adicional por espera) y el indirecto (erosión de la base de usuarios del sistema
público). Una reducción sistemática de los tiempos de espera por debajo de los
100 días no solo mejoraría la salud de los pacientes, sino que también reduciría
los incentivos a la contratación del seguro privado.

La prioridad política debería focalizarse en las comunidades autónomas con
tiempos superiores a 100 días, donde el riesgo de abandono del sistema público
es cualitativamente mayor.

\subsection{La pandemia como experimento natural}

El coeficiente positivo de $D_{\text{COVID},t}$ en el modelo de efectos aleatorios
(+17\%, significativo al 1\%) y positivo pero no significativo bajo efectos fijos
(+9\%, $p = 0{,}13$) sugiere que la pandemia, lejos de reducir el gasto
en seguros privados, lo \textbf{incrementó}. Este resultado es interpretable como
una respuesta de los hogares al colapso percibido del sistema público durante la
primera ola: ante la saturación hospitalaria y la suspensión de consultas
programadas, los hogares con capacidad de pago anticiparon la contratación de
seguros como mecanismo de protección. Este efecto es coherente con el modelo
teórico de precio sombra: la pandemia elevó drásticamente el coste implícito de
la sanidad pública, reduciendo el precio relativo del sector privado.

% ------------------------------------------------------------
\section{Contraste con la literatura previa}
% ------------------------------------------------------------

Los resultados de este trabajo coinciden en lo fundamental con la literatura
comparada. \textcite{Propper2000} obtiene elasticidades renta entre 1{,}5 y 2{,}3
para el seguro privado en Reino Unido, rango en el que se encuadra el
1{,}54 del modelo RE del presente estudio, mientras que el modelo FE preferido
arroja una elasticidad de 1{,}02, coherente con estimaciones \textit{within} que
controlan la heterogeneidad no observada. \textcite{Besley1999} documenta el efecto
de las listas de espera sobre el seguro privado en un estudio pionero para el NHS
británico, aunque sin modelizar el efecto diferido que aquí se formaliza.

La contribución específica de este trabajo es \textbf{triple}: (1) ofrece la
primera estimación sistemática con datos de panel para las CCAA españolas en el
período post-descentralización; (2) identifica el efecto umbral significativo
en los 100 días de espera como mecanismo de segmentación; (3) documenta la
persistencia del gasto en seguros privados como reflejo de la inercia en las
decisiones de consumo sanitario.
